%%
%%Huidig onboardingsproces
%%

\chapter{Huidig Onboardingsproces}
\label{ch:Huidig-Onboardingsproces}

Binnen het kader van dit onderzoek is het essentieel om eerst de huidige werkwijze in kaart te brengen. Hiervoor werd in overleg met collega’s het bestaande onboardingproces geanalyseerd. Op basis van deze analyse werden tevens de vereisten van het gemeentebestuur geïdentificeerd. Hieruit bleek dat er een duidelijke behoefte bestaat om een aanzienlijk deel van het onboardingproces binnen de ICT-dienst te automatiseren. \newline

Niet alle onderdelen van het proces lenen zich echter tot automatisatie. Zo is het selecteren van geschikte hardware, zoals laptops, en het kiezen en installeren van bijhorende software op basis van bestaande configuratiepakketten moeilijk te automatiseren. Een mogelijke oplossing zou zijn om alle toestellen in voorraad permanent operationeel te houden zodat deze automatisch geconfigureerd kunnen worden. Dit werd echter als onrealistisch beschouwd, onder meer vanwege het inefficiënte energieverbruik. \newline

Momenteel verloopt het onboardingproces grotendeels manueel. Zo worden gebruikersaccounts in Active Directory handmatig aangemaakt en worden de bijhorende groepen individueel toegewezen. Ook het toekennen van licenties en het toevoegen van gebruikers aan distributielijsten gebeurt handmatig. Daarnaast worden welkomstmails met relevante handleidingen afzonderlijk verstuurd naar nieuwe medewerkers. Eveneens wordt de welkomstbrief met het tijdelijke wachtwoord manueel afgedrukt en fysiek verwerkt in een envelop. De automatisatie van dit laatste onderdeel valt buiten de scope van dit onderzoek, aangezien dit bijkomende gespecialiseerde systemen zou vereisen. \newline

Parallel aan dit project werkt een collega aan een gelijkaardig initiatief op grotere schaal, met een focus op het automatisch versturen van meldingen naar medewerkers bij het voltooien van bepaalde stappen in het proces. Hierbij wordt gebruikgemaakt van Microsoft Power Automate. Deze keuze impliceert een zekere beperking in de selectie van alternatieve automatisatietools, omwille van compatibiliteit en integratievereisten.

Dit project komt er ook omdat er momenteel frustratie ontstaat bij werknemers omdat ze soms moeten uitzoeken of een taak al dan niet is gedaan. Zo kan het bijvoorbeeld zijn dat een collega zijn taak al heeft uitgevoerd, maar nog niet heeft teruggekoppeld. Hierdoor kan een andere collega zijn of haar werk niet doen omdat ze niet weten dat de voorwaarden om aan die stappen kunnen beginnen al gedaan zijn. Zowel om de frustraties als het een vlotter proces te maken, is dit project in leven geroepen. 

Om de werking van dat project goed te laten werken met dit project, worden er terugkoppelingen verwacht. Dit door middel van meldingen wanneer de gebruiker aangemaakt is. Andere diensten hebben immers nood aan de creatie van het mailadres om accounts in dienstspecifieke software aan te maken. Dit valt buiten de scope van de ICT dienst maar is nog steeds cruciaal voor het onboardingsproces. 

\begin{figure}[htbp]
    \centering
    \includegraphics[width=0.8\textwidth]{HuidigeOnboarding.png}
    \caption[Huidig onboardingsproces]{Deze afbeelding is het huidige onboardingsproces dat momenteel gebruikt wordt binnen de gemeente }
    \label{fig:HuidigeOnboarding}
\end{figure}
